% Independent audit and further results for research/fractional_notes.txt.
% This section uses the compensated density and Euler remainder definitions there.
% It is a fragment for the root article, not a standalone document.

\section{An exact phase diagram for fixed-total Euler comparisons}
\label{phasefrac:sec:main}

Retain the compensated density $K_{a,b}$, the scaled remainder $R_N(a,b)$,
and $w=a+b$ from the preceding fractional expansion. Let
\[
 D_{a,b}(T)=K_{a,b}(T)-K_{0,w}(T),\qquad
 q(T)=1-e^{-2T}.
\]
The exact remainder identity gives, for integers $N\ge1$,
\begin{equation}\label{phasefrac:eq:difference}
 R_N(a,b)-R_N(0,w)
 =-\int_0^\infty
 D_{a,b}(T)\frac{e^{-T}q(T)^N}{1+e^{-2T}}\,dT.
\end{equation}
The formula also defines a natural interpolation in real $N\ge1$.
All integrals below are ordinary absolutely convergent integrals after
this vanishing test factor is included. No unweighted finite-measure
assertion is made in the subcritical case.

\subsection{All truncation orders on and above the outer-order threshold}

\begin{theorem}[Axis domination and complete monotonicity]
\label{phasefrac:thm:above}
For every $a\ge1$, $b>0$, and integer $N\ge1$,
\[
 R_N(a,b)<R_N(0,a+b).
\]
Writing $B_N=R_N(0,a+b)-R_N(a,b)$ and
$\Delta B_N=B_{N+1}-B_N$, one has the stronger inequalities
\[
 (-1)^r\Delta^rB_N>0\qquad(r\ge0,\ N\ge1).
\]
In particular, $B_N$ is strictly decreasing and
$B_NB_{N+2}>B_{N+1}^2$ for every $N\ge1$.
\end{theorem}
\begin{proof}
The density formulas give the exact identity
\begin{equation}\label{phasefrac:eq:positive-density}
 D_{a,b}(T)=(I^a h_b)'(T)
    +\frac{T^{w-1}}{\Gamma(w)(e^T-1)}.
\end{equation}
For $a>1$ the first term is $I^{a-1}h_b(T)>0$, while for $a=1$
it is $h_b(T)>0$. Thus $D_{a,b}$ is strictly positive on $(0,\infty)$.
Substitution in \eqref{phasefrac:eq:difference} proves the comparison.
Taking finite differences under the integral gives
\[
 (-1)^r\Delta^rB_N
 =\int_0^\infty D_{a,b}(T)
   \frac{e^{-T}q(T)^N(1-q(T))^r}{1+e^{-2T}}\,dT>0.
\]
The same formula is a Hausdorff moment representation, with a positive
measure supported throughout $(0,1)$ after the change of variable $x=q(T)$.
Strict log-convexity follows from strict Cauchy--Schwarz for that measure.
Here $a+b>1$, so its unweighted mass is also finite; the statements
only require the moments starting at $N=1$.
\end{proof}

\begin{proposition}[The smaller scale at the threshold]
\label{phasefrac:prop:critical}
For fixed $b>0$, with $L=\tfrac12\log N$,
\[
 R_N(1,b)-R_N(0,b+1)
 \sim-\frac{L^b}{2\Gamma(b+1)N}.
\]
Thus the comparison at $a=1$ occurs on a smaller power of $N$ than
the $N^{-1/2}$ scale of the generic fractional correction.
\end{proposition}
\begin{proof}
At $a=1$, \eqref{phasefrac:eq:positive-density} becomes
\[
 D_{1,b}(T)=h_b(T)+\frac{T^b}{\Gamma(b+1)(e^T-1)}
 =\frac{e^{-T}T^b}{\Gamma(b+1)}
        \left(1+O(T^{-1})\right)
 \quad(T\to\infty).
\]
The first term $h_b(T)=O(e^{-T}T^{b-1})$ follows directly by
integrating its exponentially decaying density. Set
$t=\sqrt N e^{-T}$ in \eqref{phasefrac:eq:difference}. On the
central range, the additional $e^{-T}$ factor changes $N^{-1/2}$
to $N^{-1}$ and changes the limiting Gaussian weight to $t e^{-t^2}$.
The same dominated tail estimates used for the fractional expansion
apply, giving
\[
 R_N(0,b+1)-R_N(1,b)
 =\frac1{N\Gamma(b+1)}
   \int_{N^{-1/4}}^{N^{1/4}} t e^{-t^2}(L-\log t)^b\,dt
   +O(N^{-1}L^{b-1}),
\]
All powers here have positive arguments. Taylor expansion on this
central range and extension of the Gaussian logarithmic moments to
$(0,\infty)$ give the leading term
$L^b\int_0^\infty t e^{-t^2}\,dt=L^b/2$.
\end{proof}

\subsection{Exactly one reversal below the threshold}

The fractional integral admits a useful alternative expression. Put
\[
 \phi_a(u)=
 \begin{cases}
 1-(1-u)^{a-1},&0<u<1,\\
 1,&u>1.
 \end{cases}
\]
When $0<a<1$, this function is strictly negative below one and strictly
positive above one.

\begin{lemma}[A density with exactly one change of sign]
\label{phasefrac:lem:density}
For $0<a<1$ and $b>0$, the difference density $D_{a,b}$ has exactly
one zero $T_*=T_*(a,b)$ on $(0,\infty)$. It is positive for $T<T_*$,
negative for $T>T_*$, and $D_{a,b}'(T_*)<0$.
\end{lemma}
\begin{proof}
First derive the exact identity
\begin{equation}\label{phasefrac:eq:bose}
 \Gamma(a)\Gamma(b)T^{1-w}D_{a,b}(T)
 =\int_0^\infty\frac{\phi_a(u)u^{b-1}}{e^{Tu}-1}\,du
       +\frac{B(a,b)}{e^T-1},
 \qquad B(a,b)=\frac{\Gamma(a)\Gamma(b)}{\Gamma(w)}.
\end{equation}
Indeed, Tonelli's theorem in the definition of $I^ah_b$ gives
\[
 I^ah_b(T)=\frac1{\Gamma(a+1)\Gamma(b)}
   \int_0^\infty\frac{v^{b-1}}{e^v-1}
        \bigl(T^a-(T-v)_+^a\bigr)\,dv.
\]
Differentiation yields
\[
 (I^ah_b)'(T)=\frac1{\Gamma(a)\Gamma(b)}
 \left\{\int_0^T\frac{v^{b-1}
       [T^{a-1}-(T-v)^{a-1}]}{e^v-1}\,dv
       +T^{a-1}\int_T^\infty\frac{v^{b-1}}{e^v-1}\,dv\right\}.
\]
The cancellation in the first integral must be retained if $b\le1$.
Its integrand is $O(v^{b-1})$ at zero, and its moving-endpoint
singularity has exponent $a-1>-1$. One may justify differentiation
first with a lower cutoff in $v$, split at a fixed neighborhood of
$v=T$, and then pass to the limit; the displayed bounds give locally
uniform convergence of the derivative. Scaling $v=Tu$ and using
\eqref{phasefrac:eq:positive-density} proves \eqref{phasefrac:eq:bose}.

Let $A(T)$ denote the right side of \eqref{phasefrac:eq:bose}.
It is the transform by $K(T,u)=(e^{Tu}-1)^{-1}$ of the signed measure
\[
 d\nu(u)=\phi_a(u)u^{b-1}\,du+B(a,b)\,\delta_1(du).
\]
This measure is negative on $(0,1)$ and positive on $(1,\infty)$,
with a positive atom at the dividing point. All its transforms and
their $T$ derivatives converge absolutely on compact positive
$T$ intervals. At zero in $u$, the density is $O(u^b)$; its
singularity at $u=1$ is integrable, and the kernel controls infinity.

For $T_2>T_1>0$, the kernel ratio
\[
 r(u)=\frac{K(T_2,u)}{K(T_1,u)}
      =\frac{e^{T_1u}-1}{e^{T_2u}-1}
\]
is strictly decreasing in $u$. For example, differentiate its
logarithm and use strict increase of $x/(1-e^{-x})$ on $(0,\infty)$.
If $A(T_1)=0$, subtraction of the constant $r(1)$ gives
\[
 A(T_2)=\int_0^\infty[r(u)-r(1)]K(T_1,u)\,d\nu(u)<0.
\]
Both open intervals give strictly negative contributions, while the
atom contributes zero. Reversing this comparison shows $A(T)>0$
for every $T<T_1$. Thus $A$ can have at most one zero, and that zero
has the stated sign orientation.

The endpoint signs can also be checked explicitly. As $T\downarrow0$,
\[
 D_{a,b}(T)\sim
 \begin{cases}
 \displaystyle\frac{a}{(1-b)\Gamma(w)}T^{w-2},&0<b<1,\\[5pt]
 \displaystyle\frac{T^{a-1}\log(1/T)}{\Gamma(a)},&b=1,\\[5pt]
 \displaystyle\frac{\zeta(b)}{\Gamma(a)}T^{a-1},&b>1.
 \end{cases}
\]
For the first line, the leading term of $h_b$ is
$T^{b-1}/[(1-b)\Gamma(b)]$; differentiating its fractional integral
and adding the last term of \eqref{phasefrac:eq:positive-density}
produces the positive coefficient $a/(1-b)$, including $w=1$.
The other two lines follow from $h_1(T)\sim\log(1/T)$ and
$h_b(0)=\zeta(b)$ for $b>1$.
At infinity, the density expansion gives
\[
 D_{a,b}(T)\sim\frac{b\zeta(b+1)}{\Gamma(a-1)}T^{a-2}<0.
\]
Continuity therefore gives a zero, and the comparison proves uniqueness.

Finally, the logarithmic derivative
\[
 s_T(u)=\partial_T\log K(T,u)=-\frac{u}{1-e^{-Tu}}
\]
is strictly decreasing in $u$. At the zero,
\[
 A'(T_*)=\int[s_{T_*}(u)-s_{T_*}(1)]K(T_*,u)\,d\nu(u)<0.
\]
The positive prefactor in \eqref{phasefrac:eq:bose} proves
$D_{a,b}'(T_*)<0$.
\end{proof}

\begin{theorem}[A unique truncation-order reversal]
\label{phasefrac:thm:unique}
For every $0<a<1$ and $b>0$ there is a unique real number
$\nu(a,b)>1$ such that the interpolation in
\eqref{phasefrac:eq:difference} satisfies
\[
 R_s(a,b)-R_s(0,a+b)
 \begin{cases}
 <0,&1\le s<\nu(a,b),\\
 =0,&s=\nu(a,b),\\
 >0,&s>\nu(a,b).
 \end{cases}
\]
Its derivative with respect to $s$ is strictly positive at the zero.
Consequently the comparison for integral truncation orders reverses
exactly once, allowing equality at one integer if $\nu(a,b)$ is integral.
\end{theorem}
\begin{proof}
Write $G(s)$ for minus the interpolated difference, namely the
integral on the right of \eqref{phasefrac:eq:difference} without
the minus sign. It converges with its $s$ derivative for all $s\ge1$.
Near zero one may use
$D_{a,b}(T)=O((T^{w-2}+T^{a-1})(1+|\log T|))$.
The test $q(T)^s=O(T^s)$ makes both powers integrable because
$s+w>1$ and $s+a>0$ for $s\ge1$.
Additional logarithms from the $s$ derivatives are harmless.

For $s_2>s_1\ge1$, the ratio of the weights is
$r(T)=q(T)^{s_2-s_1}$, strictly increasing in $T$. At any zero
$G(s_1)=0$, subtracting $r(T_*)$ and using
Lemma~\ref{phasefrac:lem:density} gives
\[
 G(s_2)=\int D_{a,b}(T)
 \frac{e^{-T}q(T)^{s_1}}{1+e^{-2T}}
       [r(T)-r(T_*)]\,dT<0.
\]
The same comparison backwards gives $G(s)>0$ before that zero.
At a zero, replacing the bracket by
$\log q(T)-\log q(T_*)$ proves $G'(s)<0$.

It remains to establish existence. The manuscript's fixed-total Gaussian monotonicity theorem
\cite[\texttt{subpick:thm:fixed-total}]{baseline}, together with $E_1=0$,
gives $G(1)=R_1(0,w)-R_1(a,b)>0$. Its short comparison can be
recalled explicitly. For independent $T\sim\mathrm{Gamma}(w,1)$ and
$R\sim\mathrm{Beta}(a,b)$, the positive double kernel gives
\[
 -g_{a,b}=\mathbb E\frac{X(X+Y)}{(1+X^2)(1+Y^2)},
 \qquad X=e^{-RT},\quad Y=e^{-T}.
\]
For fixed $T>0$, the derivative with respect to $X$ is
$(2X+Y-YX^2)/((1+X^2)^2(1+Y^2))>0$.
Thus the integrand decreases strictly in $R$, and its expectation
is below the $R=0$ axis value. This proves the required initial sign.
For an independent elementary proof of the eventual negative sign,
absorb the $s=1$ weight into the signed integrable density.
Its positive part is supported below $T_*$, where $q(T)\le q(T_*)$.
Choose a fixed closed interval strictly above $T_*$; the negative
part has strictly positive mass there and $q(T)\ge q_1>q(T_*)$.
Thus, for suitable $C,c>0$,
\[
 G(s)\le Cq(T_*)^{s-1}-c q_1^{s-1}<0
\]
for sufficiently large $s$. Continuity gives a zero, and the preceding
strict comparisons prove the full statement.
\end{proof}

For every fixed $a>0$, $a\ne1$, the leading total-order powers
cancel exactly in the density difference. Its first nonzero term is
$b\zeta(b+1)T^{a-2}/\Gamma(a-1)$, with relative error $O(T^{-1})$.
The Gaussian transfer in Theorem~\ref{frac:thm:allorders} therefore gives
\begin{equation}\label{phasefrac:eq:genericscale}
 R_N(a,b)-R_N(0,a+b)
 \sim-\frac{b\zeta(b+1)\sqrt\pi}{2\Gamma(a-1)}N^{-1/2}L^{a-2}.
\end{equation}
In particular, the eventual positive side for $0<a<1$ is
\[
 R_N(a,b)-R_N(0,a+b)
 \sim-\frac{b\zeta(b+1)\sqrt\pi}{2\Gamma(a-1)}
             N^{-1/2}L^{a-2}>0\quad(0<a<1).
\]
The unique-crossing proof itself needs only the exact densities and
their sign structure, so it is logically independent of the detailed
all-order logarithmic expansion.

\begin{corollary}[A unique maximum and the derivative zero sequence]
\label{phasefrac:cor:derivatives}
For $0<a<1$ and $b>0$, let
$\mathcal D(s)=R_s(a,b)-R_s(0,a+b)$ denote the interpolation above.
For every $r\ge0$, the derivative $\mathcal D^{(r)}$ has exactly one
zero $\nu_r$ on $[1,\infty)$. These zeros are simple and satisfy
\[
 1<\nu_0<\nu_1<\nu_2<\cdots,\qquad \nu_r\longrightarrow\infty.
\]
The sign of $\mathcal D^{(r)}$ before its zero is $(-1)^{r+1}$,
and its sign afterwards is $(-1)^r$. In particular, $\mathcal D$
increases to a unique positive maximum at $\nu_1$ and then decreases
to zero.
\end{corollary}
\begin{proof}
Set $\ell(T)=-\log q(T)>0$ and
\[
 H_r(s)=\int_0^\infty D_{a,b}(T)
 \frac{e^{-T}q(T)^s\ell(T)^r}{1+e^{-2T}}\,dT.
\]
Then $\mathcal D^{(r)}=(-1)^{r+1}H_r$ and $H_r'=-H_{r+1}$.
Every $H_r$ has the same single-change density signs as $D_{a,b}$.
The preceding ratio argument shows that it has at most one zero,
with strictly negative derivative there, and support separation shows
that it is negative for all sufficiently large $s$.

Theorem~\ref{phasefrac:thm:unique} supplies the first zero $\nu_0$.
At any zero of $H_r$, one has
$H_{r+1}(\nu_r)=-H_r'(\nu_r)>0$. Thus $H_{r+1}$ has exactly one
zero $\nu_{r+1}>\nu_r$. This proves induction and all sign assertions.
To see that the zero sequence diverges, fix $s\ge1$. The negative
part of the integrand is supported where $\ell(T)\le\ell(T_*)$.
Choose a positive-mass closed interval below $T_*$ on which
$\ell(T)\ge\ell_1>\ell(T_*)$. Its positive contribution dominates
the entire negative part as $r\to\infty$. Hence $H_r(s)>0$ for
large $r$, so $s<\nu_r$. Finally $\mathcal D(s)\to0$ by dominated
convergence using the $s=1$ absolute integrable weight.
\end{proof}

\paragraph{A precise next asymptotic question.}
For fixed $b>0$, the unique crossing order tends to infinity as
$a\uparrow1$: continuity at each fixed $s$ and the strict comparison
at $a=1$ imply this directly. Balancing the generic fractional
correction against the critical $N^{-1}$ term suggests a boundary
regime for $\nu(a,b)$, but proving that regime requires an expansion
uniform in $1-a$ and $N$. It is not asserted here.

\subsection{An exact discrete transition at half orders}

\begin{corollary}[The first positive integer comparison is $21$]
\label{phasefrac:cor:twentyone}
Let $\Delta_N=R_N(1/2,1/2)-R_N(0,1)$. Then
\[
 \Delta_N<0\quad(1\le N\le20),\qquad
 \Delta_N>0\quad(N\ge21).
\]
In the continuous interpolation,
$20<\nu(1/2,1/2)<21$.
\end{corollary}
\begin{proof}
The exact rational interval calculation supplied with the article gives
\begin{align*}
 -0.000325283275869852646927
 &<\Delta_{20}<-0.000325283275869852646926,\\
  \phantom{-}0.000023323498144449794846
 &<\Delta_{21}<\phantom{-}0.000023323498144449794847.
\end{align*}
Theorem~\ref{phasefrac:thm:unique} then proves the whole integer
classification.
\end{proof}

Here are the details of the finite certificate. For each $N\in\{20,21\}$,
put $M=N+100$ and compute
\[
 2^N\{E_N(1/2,1/2)-E_M(1/2,1/2)
       -E_N(0,1)+E_M(0,1)\}.
\]
The inherited bound $0<R_M<57/50$~\cite{incoming-sharp} makes its
distance from $\Delta_N$ less than $(57/50)2^{N-M}$. This is one
radius: the difference of two numbers in $(0,57/50)$ has magnitude
less than $57/50$. The direct finite Euler weights are
\[
 E_N=\sum_{n=0}^{N-1}
 \frac{(-1)^n}{2^N}\left(\sum_{j=n+1}^{N}\binom Nj\right)f_n.
\]
For every positive integer $r$, integer square roots produce rational
dyadic bounds on $r^{-1/2}$:
\[
 k=\left\lfloor\sqrt{\left\lfloor 2^{2p}/r\right\rfloor}\right\rfloor,
 \qquad
 \frac{k}{2^p}\le r^{-1/2}<\frac{k+1}{2^p},
 \qquad p=N+200.
\]
All harmonic sums, products, Euler weights, and tail additions then
use rational interval arithmetic. No floating-point operation
contributes to acceptance of the signs.

\begin{table}[tbp]
\centering\small
\begin{tabular}{@{}rr@{}}
\toprule
$N$ & Outward rounded enclosure for $\Delta_N$\\
\midrule
1 & $(-0.324097778244,\,-0.324097778243)$\\
8 & $(-0.015731327706,\,-0.015731327705)$\\
16 & $(-0.002351554584,\,-0.002351554583)$\\
20 & $(-0.000325283276,\,-0.000325283275)$\\
21 & $(0.000023323498,\,0.000023323499)$\\
32 & $(0.001999520010,\,0.001999520011)$\\
64 & $(0.002856035510,\,0.002856035511)$\\
128 & $(0.002540045225,\,0.002540045226)$\\
256 & $(0.001952284331,\,0.001952284332)$\\
\bottomrule
\end{tabular}
\caption{Exact rational enclosures for the fixed-total half-order
comparison. The certificate files retain all rational endpoints and
24 outward rounded decimal places. Only the signs at $20$ and $21$
are needed for Corollary~\ref{phasefrac:cor:twentyone}; the remaining
rows illustrate the subsequent positive maximum and decay.}
\label{phasefrac:tab:reversal}
\end{table}

The threshold replay imports the adjacent general fractional verifier.
Its dependency and source hash are recorded explicitly; it is not a
second arithmetic implementation. The global sign theorem and the
finite enclosure computation provide different parts of the proof.

